% Options for packages loaded elsewhere
\PassOptionsToPackage{unicode}{hyperref}
\PassOptionsToPackage{hyphens}{url}
\documentclass[
]{article}
\usepackage{xcolor}
\usepackage[margin=1in]{geometry}
\usepackage{amsmath,amssymb}
\setcounter{secnumdepth}{-\maxdimen} % remove section numbering
\usepackage{iftex}
\ifPDFTeX
  \usepackage[T1]{fontenc}
  \usepackage[utf8]{inputenc}
  \usepackage{textcomp} % provide euro and other symbols
\else % if luatex or xetex
  \usepackage{unicode-math} % this also loads fontspec
  \defaultfontfeatures{Scale=MatchLowercase}
  \defaultfontfeatures[\rmfamily]{Ligatures=TeX,Scale=1}
\fi
\usepackage{lmodern}
\ifPDFTeX\else
  % xetex/luatex font selection
\fi
% Use upquote if available, for straight quotes in verbatim environments
\IfFileExists{upquote.sty}{\usepackage{upquote}}{}
\IfFileExists{microtype.sty}{% use microtype if available
  \usepackage[]{microtype}
  \UseMicrotypeSet[protrusion]{basicmath} % disable protrusion for tt fonts
}{}
\makeatletter
\@ifundefined{KOMAClassName}{% if non-KOMA class
  \IfFileExists{parskip.sty}{%
    \usepackage{parskip}
  }{% else
    \setlength{\parindent}{0pt}
    \setlength{\parskip}{6pt plus 2pt minus 1pt}}
}{% if KOMA class
  \KOMAoptions{parskip=half}}
\makeatother
\usepackage{color}
\usepackage{fancyvrb}
\newcommand{\VerbBar}{|}
\newcommand{\VERB}{\Verb[commandchars=\\\{\}]}
\DefineVerbatimEnvironment{Highlighting}{Verbatim}{commandchars=\\\{\}}
% Add ',fontsize=\small' for more characters per line
\usepackage{framed}
\definecolor{shadecolor}{RGB}{248,248,248}
\newenvironment{Shaded}{\begin{snugshade}}{\end{snugshade}}
\newcommand{\AlertTok}[1]{\textcolor[rgb]{0.94,0.16,0.16}{#1}}
\newcommand{\AnnotationTok}[1]{\textcolor[rgb]{0.56,0.35,0.01}{\textbf{\textit{#1}}}}
\newcommand{\AttributeTok}[1]{\textcolor[rgb]{0.13,0.29,0.53}{#1}}
\newcommand{\BaseNTok}[1]{\textcolor[rgb]{0.00,0.00,0.81}{#1}}
\newcommand{\BuiltInTok}[1]{#1}
\newcommand{\CharTok}[1]{\textcolor[rgb]{0.31,0.60,0.02}{#1}}
\newcommand{\CommentTok}[1]{\textcolor[rgb]{0.56,0.35,0.01}{\textit{#1}}}
\newcommand{\CommentVarTok}[1]{\textcolor[rgb]{0.56,0.35,0.01}{\textbf{\textit{#1}}}}
\newcommand{\ConstantTok}[1]{\textcolor[rgb]{0.56,0.35,0.01}{#1}}
\newcommand{\ControlFlowTok}[1]{\textcolor[rgb]{0.13,0.29,0.53}{\textbf{#1}}}
\newcommand{\DataTypeTok}[1]{\textcolor[rgb]{0.13,0.29,0.53}{#1}}
\newcommand{\DecValTok}[1]{\textcolor[rgb]{0.00,0.00,0.81}{#1}}
\newcommand{\DocumentationTok}[1]{\textcolor[rgb]{0.56,0.35,0.01}{\textbf{\textit{#1}}}}
\newcommand{\ErrorTok}[1]{\textcolor[rgb]{0.64,0.00,0.00}{\textbf{#1}}}
\newcommand{\ExtensionTok}[1]{#1}
\newcommand{\FloatTok}[1]{\textcolor[rgb]{0.00,0.00,0.81}{#1}}
\newcommand{\FunctionTok}[1]{\textcolor[rgb]{0.13,0.29,0.53}{\textbf{#1}}}
\newcommand{\ImportTok}[1]{#1}
\newcommand{\InformationTok}[1]{\textcolor[rgb]{0.56,0.35,0.01}{\textbf{\textit{#1}}}}
\newcommand{\KeywordTok}[1]{\textcolor[rgb]{0.13,0.29,0.53}{\textbf{#1}}}
\newcommand{\NormalTok}[1]{#1}
\newcommand{\OperatorTok}[1]{\textcolor[rgb]{0.81,0.36,0.00}{\textbf{#1}}}
\newcommand{\OtherTok}[1]{\textcolor[rgb]{0.56,0.35,0.01}{#1}}
\newcommand{\PreprocessorTok}[1]{\textcolor[rgb]{0.56,0.35,0.01}{\textit{#1}}}
\newcommand{\RegionMarkerTok}[1]{#1}
\newcommand{\SpecialCharTok}[1]{\textcolor[rgb]{0.81,0.36,0.00}{\textbf{#1}}}
\newcommand{\SpecialStringTok}[1]{\textcolor[rgb]{0.31,0.60,0.02}{#1}}
\newcommand{\StringTok}[1]{\textcolor[rgb]{0.31,0.60,0.02}{#1}}
\newcommand{\VariableTok}[1]{\textcolor[rgb]{0.00,0.00,0.00}{#1}}
\newcommand{\VerbatimStringTok}[1]{\textcolor[rgb]{0.31,0.60,0.02}{#1}}
\newcommand{\WarningTok}[1]{\textcolor[rgb]{0.56,0.35,0.01}{\textbf{\textit{#1}}}}
\usepackage{graphicx}
\makeatletter
\newsavebox\pandoc@box
\newcommand*\pandocbounded[1]{% scales image to fit in text height/width
  \sbox\pandoc@box{#1}%
  \Gscale@div\@tempa{\textheight}{\dimexpr\ht\pandoc@box+\dp\pandoc@box\relax}%
  \Gscale@div\@tempb{\linewidth}{\wd\pandoc@box}%
  \ifdim\@tempb\p@<\@tempa\p@\let\@tempa\@tempb\fi% select the smaller of both
  \ifdim\@tempa\p@<\p@\scalebox{\@tempa}{\usebox\pandoc@box}%
  \else\usebox{\pandoc@box}%
  \fi%
}
% Set default figure placement to htbp
\def\fps@figure{htbp}
\makeatother
\setlength{\emergencystretch}{3em} % prevent overfull lines
\providecommand{\tightlist}{%
  \setlength{\itemsep}{0pt}\setlength{\parskip}{0pt}}
\usepackage{bookmark}
\IfFileExists{xurl.sty}{\usepackage{xurl}}{} % add URL line breaks if available
\urlstyle{same}
\hypersetup{
  pdftitle={Bayesian Hierarchical Models},
  hidelinks,
  pdfcreator={LaTeX via pandoc}}

\title{Bayesian Hierarchical Models}
\author{}
\date{\vspace{-2.5em}2026-04-17}

\begin{document}
\maketitle

\subsection{Introduction}\label{introduction}

A hierarchical Bayesian model embeds the group-level parameters of a
multilevel design inside a higher-level prior, so that the group-level
estimates inform each other rather than being computed in isolation. The
pay-off is partial pooling: groups with abundant data are estimated
almost from their own observations, groups with little data are pulled
toward the overall mean, and every estimate carries a credible interval
that reflects both the within-group sample size and the between-group
heterogeneity. This single mechanism resolves a long list of practical
problems --- small-sample variance, unbalanced designs, sparse counts
--- that frequentist mixed-effects models can only address by
reweighting or by appealing to asymptotic theory.

\subsection{Prerequisites}\label{prerequisites}

A working understanding of Bayesian inference, the random-effects mixed
model, and the role of exchangeability in motivating shared priors.

\subsection{Theory}\label{theory}

A two-level model has

\begin{itemize}
\tightlist
\item
  Level 1: \(y_{ij} \mid \theta_i \sim p(y \mid \theta_i)\) --- the data
  given group-specific parameters.
\item
  Level 2: \(\theta_i \mid \mu, \tau \sim p(\theta \mid \mu, \tau)\) ---
  group-specific parameters drawn from a population distribution.
\item
  Level 3 (hyperprior): \(\mu \sim p(\mu)\), \(\tau \sim p(\tau)\) ---
  priors on the population parameters.
\end{itemize}

The posterior of \(\theta_i\) is a precision-weighted average of the
within-group estimate (driven by the data in group \(i\)) and the
population mean \(\mu\), with weight inversely proportional to the
within-group variance and to the cross-group variance \(\tau^2\).
Smaller groups rely more on the population mean; larger groups stand on
their own data. This is the formal expression of ``shrinkage'' or
``borrowing strength.''

\subsection{Assumptions}\label{assumptions}

Exchangeability of groups (group labels carry no information beyond what
is in the predictors) and a hyperprior whose support spans plausible
values of \(\mu\) and \(\tau\). A poorly chosen hyperprior ---
particularly an overly tight prior on \(\tau\) --- can dominate the
posterior in small samples and look like a strong empirical claim.

\subsection{R Implementation}\label{r-implementation}

\begin{Shaded}
\begin{Highlighting}[]
\FunctionTok{library}\NormalTok{(brms)}

\CommentTok{\# 30 hospitals, events and trials per hospital}
\FunctionTok{set.seed}\NormalTok{(}\DecValTok{2026}\NormalTok{)}
\NormalTok{n\_hosp }\OtherTok{\textless{}{-}} \DecValTok{30}
\NormalTok{n\_per  }\OtherTok{\textless{}{-}} \FunctionTok{sample}\NormalTok{(}\DecValTok{5}\SpecialCharTok{:}\DecValTok{50}\NormalTok{, n\_hosp, }\AttributeTok{replace =} \ConstantTok{TRUE}\NormalTok{)}
\NormalTok{true\_rates }\OtherTok{\textless{}{-}} \FunctionTok{rbeta}\NormalTok{(n\_hosp, }\DecValTok{3}\NormalTok{, }\DecValTok{12}\NormalTok{)}
\NormalTok{events }\OtherTok{\textless{}{-}} \FunctionTok{rbinom}\NormalTok{(n\_hosp, n\_per, true\_rates)}

\NormalTok{d }\OtherTok{\textless{}{-}} \FunctionTok{data.frame}\NormalTok{(}\AttributeTok{hospital =} \FunctionTok{factor}\NormalTok{(}\DecValTok{1}\SpecialCharTok{:}\NormalTok{n\_hosp),}
                \AttributeTok{events =}\NormalTok{ events, }\AttributeTok{trials =}\NormalTok{ n\_per)}

\NormalTok{fit }\OtherTok{\textless{}{-}} \FunctionTok{brm}\NormalTok{(events }\SpecialCharTok{|} \FunctionTok{trials}\NormalTok{(trials) }\SpecialCharTok{\textasciitilde{}} \DecValTok{1} \SpecialCharTok{+}\NormalTok{ (}\DecValTok{1} \SpecialCharTok{|}\NormalTok{ hospital),}
           \AttributeTok{data =}\NormalTok{ d, }\AttributeTok{family =}\NormalTok{ binomial,}
           \AttributeTok{chains =} \DecValTok{4}\NormalTok{, }\AttributeTok{iter =} \DecValTok{2000}\NormalTok{, }\AttributeTok{refresh =} \DecValTok{0}\NormalTok{)}
\FunctionTok{summary}\NormalTok{(fit)}
\end{Highlighting}
\end{Shaded}

\subsection{Output \& Results}\label{output-results}

\texttt{summary(fit)} returns the population-level intercept (overall
log-odds), the between-hospital SD on the logit scale, and per-hospital
posterior summaries via \texttt{ranef(fit)}. The shrinkage of
small-hospital estimates toward the population mean is most visible when
you plot raw event proportions against posterior means alongside
hospital sample size.

\subsection{Interpretation}\label{interpretation}

A reporting sentence: ``A Bayesian binomial hierarchical model estimated
the population baseline at \(\mathrm{logit}^{-1}(-1.4) = 0.20\) (95 \%
credible interval 0.16 to 0.24); the between-hospital standard deviation
on the logit scale was 0.51 (95 \% CrI 0.34 to 0.74), and small-hospital
posterior means were drawn substantially toward the population
baseline.'' Always report cluster-specific and overall estimates; one
without the other misses the central insight of the hierarchical fit.

\subsection{Practical Tips}\label{practical-tips}

\begin{itemize}
\tightlist
\item
  Partial pooling is the headline benefit: do not throw it away by
  switching to a fixed-effects fit just because cluster counts are
  large.
\item
  For divergent transitions, switch to non-centred parameterisation;
  \texttt{brms} does this automatically when you specify
  \texttt{prior(...\ ,\ class\ =\ "sd",\ ...)}.
\item
  The prior on \(\tau\) matters enormously when the number of clusters
  is small; report a sensitivity analysis with at least one wider and
  one tighter prior.
\item
  For more than two levels (e.g.~patient \(\subset\) hospital
  \(\subset\) region), nest random effects:
  \texttt{(1\ \textbar{}\ region/hospital)}. The same shrinkage logic
  applies at each level.
\item
  Do not interpret cluster-specific posterior means as if they came from
  independent fits; they are \emph{intentionally} shrunk and that
  shrinkage is the right answer to ``how good is hospital \(i\)?''
\end{itemize}

\subsection{R Packages Used}\label{r-packages-used}

\texttt{brms} for the hierarchical Bayesian fit, \texttt{rstanarm} for
pre-compiled \texttt{stan\_glmer} alternatives, \texttt{tidybayes} for
tidy cluster-level draws when bespoke shrinkage plots are needed, and
\texttt{bayesplot} for diagnostic visualisation.

\subsection{For Reviewers}\label{for-reviewers}

\textbf{What to look for in a paper using this method.}

\begin{itemize}
\tightlist
\item
  \textbf{Common misapplications.}
\item
  \textbf{Diagnostics that should be reported but often aren't.}
\item
  \textbf{Red flags in tables and figures.}
\item
  \textbf{What to verify.}
\item
  \textbf{What an adequate Methods paragraph must contain.}
\end{itemize}

\subsection{See also --- labs in this
chapter}\label{see-also-labs-in-this-chapter}

\begin{itemize}
\tightlist
\item
  \href{labs/lab-bayesian-thinking-control-vs-decision-error.Rmd}{Bayesian
  thinking; α control vs decision error}
\item
  \href{labs/lab-brms-stan-loo-hierarchical-models.Rmd}{brms / Stan,
  LOO, hierarchical models}
\end{itemize}

\end{document}
